\scanpage{12}
Or $Q_i$ n'a \deleted{\unsure{aucun}}\inserted{pas de} terme ne contenant aucun des arguments après la virgule;
en plus, notons que $(\eta^r)^n=0$ si $n>1$. Donc on trouve
\[
Q_i(\text{---})=\begin{cases}
0&\text{si }i<r\\
(-1)^{r-1}(r-1)!\,Q_{i-r}^{(r)}(c^1(x),\text{---}\,c^i(x))
\end{cases}
\]
où $Q_{i-r}^{(r)}(c^1\deleted{\unsure{(x)}},\text{---},c^i\deleted{\unsure{(x)}})$
est le coefficient de $d^r$ dans $Q_i(c^1,\text{---}\,c^i;d^1,\text{---}\,d^i)$.
$(r\leq i)$.
\inserted{(est un \unsure{pol. isob.} de poids $i-r$)}
Ainsi on trouve, tenant compte que $(\eta^r)^n=0$ pour $n>1$
\[
\widetilde C(i_!(x))=1+(-1)^{r-1}(r-1)!\eta^r
\left[\varepsilon(x)+\sum_{i=0}^\infty Q_i^{(r)}(c^1(x),\text{---}\,c^i(x))\right]
\]
\unsure{Comparons|Convenons} \unsure{avec|pour} R.R.
\[
\widetilde C(i_!(x))=\deleted{\unsure{\widetilde C}}(i_*)_r(\widetilde C(x))
\]
\unsure{on écrit} \unsure{que}
\[
(i_*)_r\left[n,\sum_0^\infty x^i\right]
=1+(-1)^{r-1}(r-1)!\eta^r\biggl[n+\sum_{i=0}^\infty Q_i^{(r)}(x^1\text{---}x^i)
\]

N.B. $Q_0^{(r)}(c^1,\text{---}\,c^r)=0$, i.e.
$Q_r^{\unsure{*}}(c^1\text{---}c^r;0,\text{---}\,0,d^r)=0$

i.e. $(1+c^1+\text{---}+c^r)*(1+d^r)$ est de filtration $\geq r+1$

\unsure{De m.}, on voit
\[
\uuline{Q_1^{(r)}(c^1)}\deleted{\unsure{c^2}}
=-\frac{(r+1-1)!}{(r-1)!(1-1)!}\,c^1=\uuline{-rc^1}
\]

On fera attention que, par suite, s'il y a de la torsion dans $A(X)$,
l'hom $(i_*)_r$ est non injectif, à cause de \unsure{ses|les} coefficients numériques
(\unsure{dus tant} à $(r-1)!$ \unsure{qu'à ceux} provenant des $Q_i^{(r)}$

\margin{11}
